% RESEARCH_DRAFT_SHA256: 52482d92a001f4de5abc47975db8627da26c993856df40302d2166107d31fb97
% Layout and source-reference wrappers expanded for the web renderer.
% FULL_SOLUTION_SCOPE: {"target": "Entire cyclic sumset recognition problem for sparse binary residue lists and explicit membership arrays, as in the original scope note.", "result": "NP-complete under either encoding, even for odd moduli. A polynomial-time algorithm exists if and only if P=NP.", "coverage": "All input sets and moduli covered by NP membership; polynomial reductions include empty sets, odd integer minima, no-wrap proof and dense encoding size. No weakening of original question or source metadata.", "dependencies": "External integer NP-completeness Theorem1.2 of Abboud, Fischer, Safier and Wallheimer, arxiv2410.18661v2, including polynomial-interval reduction size. Its gadget proof is not independently re-proved here.", "novelty": "No novelty asserted. Root Codex GPT-6 supplied the elementary cyclic reduction; Astra independently inspected and wrote it.", "human_verification": "None. Codex GPT-6 written proof review and Astra self-audit only.", "unverified": "No proof of P not equal NP; no independent verification of the external integer gadgets; no exhaustive priority search."}
\documentclass[11pt]{article}
\usepackage[margin=25mm]{geometry}
\usepackage{fontspec,unicode-math}
\setmainfont{Latin Modern Roman}
\setmathfont{Latin Modern Math}
\usepackage{amsmath,amssymb,mathtools}
\usepackage{xurl,hyperref}
\hypersetup{colorlinks=true,urlcolor=blue,linkcolor=blue}
\setcounter{secnumdepth}{0}
\setlength{\parindent}{0pt}
\setlength{\parskip}{5pt}
\setlength{\emergencystretch}{3em}
\newcommand{\sref}[2]{\hyperlink{source:#1:#2}{[S#2]}}
\newcommand{\cataloguescope}[1]{\paragraph{Source question.}#1}
\newcommand{\floor}[1]{\left\lfloor#1\right\rfloor}
\begin{document}
\section*{Research-attempt provenance}
This standalone research attempt retains the complete original definition body below,
including its historical source and preparation notes. Those notes describe the
earlier catalogue preparation; the new attempt and its verification are identified
separately after that body. No original ID, catalogue status, or source statement is changed.

\section{Efficient Recognition of Cyclic Sumsets}\label{statement}

\subsection{Definitions and mathematical statement}

\paragraph{Shared notation.}Standard number systems are denoted by $\mathbb N,\mathbb Z,\mathbb Q,\mathbb R,\mathbb C$; index ranges, computational models and specialized notation follow the cited problem source and the definitions below.


The cyclic group $\mathbb Z/n\mathbb Z$ uses addition modulo $n$. For a subset $A$, its sumset is $A+A=\{a+a':a,a'\in A\}$; summands need not be distinct. The decision input consists of $n$ and an encoding of $B\subseteq\mathbb Z/n\mathbb Z$. A positive instance has a subset $A$ whose sumset equals $B$ exactly, rather than merely containing it.
\paragraph{Statement recorded in the supplied source.}
Given $B\subseteq\mathbb Z/n\mathbb Z$, is there an efficient algorithm that decides whether some $A\subseteq\mathbb Z/n\mathbb Z$ satisfies $B=A+A$? \href{https://doi.org/10.4171/owr/2022/22}{[S2]}
\paragraph{Scope qualification.}
The complexity depends on the source encoding of $B$ and $n$; explicit membership arrays and sparse binary-encoded lists have different input lengths. No encoding-independent polynomial-time claim is inserted.
\paragraph{Source status.}
The supplied catalogue labels this record open. Its source-status metadata was checked against the supplied snapshot on 2026-10-11. This import does not independently establish current open status or mathematical correctness. The catalogue literature assessment is dated 2026-08-22.
\subsection{Short English statement}
Can one efficiently recognize which subsets of a finite cyclic group are sums of a subset with itself?
\subsection{Sources}
\begin{itemize}
\item[S1] Ulam AI and the UnsolvedMath Contributors, UnsolvedMath: A Curated Collection of Open Mathematics Problems, catalogue entry. Catalogue curation and metadata: CC BY 4.0. The supplied source question is retained; definitions and the English summary are editorial additions.
\url{https://huggingface.co/datasets/ulamai/UnsolvedMath}.
\item[S2] Combinatorics, Probability and Computing (2023), Oberwolfach Reports, DOI 10.4171/owr/2022/22, pp. 1165--1237
\url{https://doi.org/10.4171/owr/2022/22}.
\item[S3] official OWR report
\url{https://publications.mfo.de/bitstream/handle/mfo/3964/OWR_2022_22.pdf?isAllowed=y&sequence=4}.
\textit{Underlying Oberwolfach Reports material retains its source-specific licence, including CC BY-SA 4.0 where applicable; see THIRD\_PARTY\_NOTICES.md.}
\end{itemize}
\paragraph{Preparation and review.}
Original question attributed to the cited source and UnsolvedMath contributors. Codex (GPT-6; exact serving revision not exposed) assisted with definitions, English summaries, source inspection and formatting. No new human mathematical authorship or independent proof verification is asserted. Catalogue identity checked against the complete supplied registry and existing repository statements; known flagged repetitions and obvious identical conjectures excluded. Full mathematical-equivalence/literature review has not been performed. Statement preparation and final review followed contribution prompts C and D; automated checks verify structure and identity, not open status.
\label{end:statement}

\section{Literature and context check}

I read Problem 15, attributed to Alon and Granville and presented by Marcelo Campos,
on printed page 1228 of the primary report. It asks exactly the cyclic recognition
question preserved above. The sumset-counting talk on pages 1174--1175 is a different
passage and is not the location of this problem.
I also consulted the introduction, Theorem 1.2 and the accompanying reduction-size
paragraph of Abboud, Fischer, Safier and Wallheimer, \emph{Recognizing Sumsets is
NP-Complete}, arXiv:2410.18661v2. Their theorem states NP-completeness over the
integers; the accompanying paragraph gives a reduction from 3-SAT with $N$
variables and clauses to a set contained in an interval of size $O(N^4)$.
That integer hardness theorem is an explicitly external dependency: its positioning
and masking gadgets are not re-proved or independently verified here.
The elementary reduction below was suggested by the root Codex GPT-6 agent and
then independently inspected, written out and adversarially checked by this
OpenAI gpt-6-astra agent. No novelty claim is made.

\section{Research attempt: Integer-to-cyclic recognition reduction with encoding control}

\paragraph{Decision problems and encodings.}
Let $q\geq1$ be an integer and let $B\subseteq\mathbb Z/q\mathbb Z$.
The question is whether $B=A+A$ for some subset $A$, with addition modulo $q$.
Consider either (i) $q$ in binary and the distinct canonical residues of $B$ in a
binary list, or (ii) a membership array of length $q$. We prove NP-completeness
under either encoding, using the stated external integer theorem. Thus a
polynomial-time decision algorithm exists if and only if $\mathrm P=\mathrm{NP}$.
This conclusion does not assert an unconditional separation of those classes.

\paragraph{Membership in NP.}
The empty input set has the empty root. For a nonempty root choose $a_0\in A$.
The injection $a\mapsto a+a_0$ maps $A$ into $B$, so $|A|\leq |B|$.
A witness list therefore uses $O(|B|\log(q+1))$ bits. Computing all at most
$|B|^2$ pair sums, reducing modulo $q$, removing repetitions and comparing with
$B$ verifies the witness in polynomial time in either specified input length.
This also covers $q=1$; there is no hidden demand for a witness of size $q$ under
sparse encoding.

\paragraph{Normalization of integer instances.}
Given a finite nonempty set $S\subseteq\mathbb Z$, put $s_0=\min S$.
If $S=C+C$, then $s_0=2\min C$. Thus odd $s_0$ implies a NO instance.
For even $s_0$ define $B=S-s_0$ and $M=\max B$. Then
\[
 S=C+C\quad\Longleftrightarrow\quad
 B=(C-s_0/2)+(C-s_0/2).
\]
Conversely any integer root $D$ of $B$ yields the root $D+s_0/2$ of $S$.
Here $0\in B\subseteq[0,M]\cap\mathbb Z$ and $M\geq0$.
An empty $S$ is mapped to a fixed YES instance, for example $(q,B)=(3,\varnothing)$.
An odd minimum is mapped to the fixed NO instance $(3,\{0,1\})$.
To verify the latter, a root has at most two elements by the injection just proved.
A singleton has a singleton sumset. For two distinct residues $a,b$ modulo $3$,
the three sums $2a,a+b,2b$ are distinct: equality of either adjacent pair implies
$a=b$, and $2a=2b$ does also because $2$ is invertible modulo $3$.
Hence its sumset has three elements.

\paragraph{The no-wrap lemma.}
For a normalized input set use the odd modulus
\[
 q=2M+3.
\]
We claim that $B$ has an integer sumset root if and only if it has a root modulo
$q$. The forward implication follows just by reducing an integer equality
modulo $q$; distinct elements of $B\subseteq[0,M]$ remain distinct.
For the reverse implication suppose $B=A+A$ modulo $q$. Because $0\in B$,
choose $a,b\in A$ with $a+b=0$ modulo $q$.
Let $r,s\in[0,q-1]$ be the canonical residues of $2a,2b$. Both belong to $B$,
so $0\leq r,s\leq M$. Moreover $r+s=0$ modulo $q$ and
$0\leq r+s\leq2M<q$. Therefore $r=s=0$.
Since $q$ is odd, multiplication by $2$ is injective, and $a=b=0$ modulo $q$.
We have proved $0\in A$. It follows that every canonical representative of $A$
is in $B$, using $a=0+a\in A+A$. All representatives of $A$ thus lie in
$[0,M]$. Their ordinary pair sums lie in $[0,2M]\subset[0,q-1]$, so none wraps
modulo $q$. The modular equality is consequently the ordinary integer equality
$A+A=B$. This proves both directions.

\paragraph{Polynomial encoding bounds.}
Subtracting $s_0$ and forming $q=2M+3$ increases the binary length of each
integer by at most an additive constant relative to the largest input magnitude.
The number of elements of $B$ is $|S|$ and the exceptional instances have
constant size. The map is thus a polynomial-time many-one reduction for sparse
binary lists. For membership arrays, arbitrary huge binary integer inputs would
not give a polynomial-size output array. Instead compose directly with the
external 3-SAT reduction whose integer interval has size $O(N^4)$.
After normalization $M=O(N^4)$ and hence $q=O(N^4)$, so the whole membership
array can be produced in polynomial time. The same theorem consequently proves
NP-hardness for both encodings. Together with membership in NP this establishes
the claimed classification, even with hardness restricted to odd $q$.

\paragraph{An exact algorithm and the failed shortcut.}
For completeness there is an elementary exponential search with candidate set
\[
 H=\{a\in\mathbb Z/q\mathbb Z:2a\in B\}.
\]
Every root lies in $H$, since every self-sum must belong to $B$.
For each $b\in B$, the congruence $2a=b\pmod q$ has at most two solutions,
computable by the Euclidean algorithm (one for odd $q$). Thus $|H|\leq2|B|$.
Enumerating subsets of $H$ and checking pair sums is a finite exact algorithm.
The necessary condition $A\subseteq H$ does not itself identify the root, and
checking $H+H=B$ is not a correct recognition algorithm: for $q=7$ and
$B=\{0,1,2\}$ the root $A=\{0,1\}$ works, but $H=\{0,1,4\}$ contains
$4$ and $H+H$ contains $4\notin B$.
The hardness classification concerns arbitrary inputs, not the termination of
this exponential algorithm or the possibility of useful special-case algorithms.

\section{Verification and adversarial audit}

The reverse implication was checked at its two essential points: $0\in B$ alone
does not imply $0\in A$ in a general cyclic group, and an arbitrarily large
modulus alone does not prevent a modular root from having large representatives.
The doubled-residue argument proves both claims using oddness and $2M<q$.
For example the use of oddness cannot simply be erased: modulo $2$, the set
$A=\{1\}$ has $A+A=\{0\}$ without containing $0$.
The cases $S=\varnothing$, $M=0$, odd negative minima, $q=1$ in the verification
problem, singleton roots and sparse witness size were separated explicitly.
The dense encoding proof uses the external polynomial-interval theorem, not
an unjustified expansion of arbitrary binary integers into an array.

\paragraph{Audit corrections.} The source location was corrected from the separate Campos counting talk to
Problem 15 on page 1228. Oddness, empty inputs, negative-minimum normalization,
and polynomial array size were added explicitly. The false maximal-candidate-set
recognition shortcut was falsified by the displayed modulo-$7$ example.
\paragraph{External computation.} None; symbolic arguments and boundary checks only.
\section{Scope and remaining gaps}

The argument supplies the usual complexity-theoretic classification of the entire
cyclic recognition question for both explicit encodings identified in the preserved
scope note. It depends on the cited integer NP-completeness theorem and on the
standard definitions of polynomial-time reduction and NP-completeness. It neither
proves $\mathrm P\ne\mathrm{NP}$ nor settles fine-grained, average-case, or special
input-family questions that the original statement did not ask.
\paragraph{Full original-target declaration.}
The entire original decision question is classified for both stated encodings,
without restricting the input sets, weakening the requested decision, or replacing
cyclic groups by a different family. The elementary reduction and both encoding
arguments are fully proved here. The integer hardness theorem is a cited external
dependency. The parent Codex GPT-6 agent has reviewed the complete written
reduction and approved this full scope. No claim of an unconditional
$\mathrm P\ne\mathrm{NP}$ theorem or new mathematical priority is made.

\paragraph{Integration review.} Codex (GPT-6; exact serving revision not exposed) reviewed the written arguments and expanded layout and source-reference wrappers. Repository metadata and the production web renderer were checked using Python 3.13 and JavaScript with MathJax 3.2.2. This model-assisted review is not independent human proof verification; no PDF compilation or formal proof checking is asserted.

\section{Completion Estimate}
\noindent\textbf{COMPLETION: 100\%}
The proof gives the complete NP-completeness classification for both specified
encodings of the original recognition problem. Completion measures original-target
coverage, not novelty or independent verification.
\section{References}
The original source references and catalogue attribution are preserved in the
Sources section above. The following URLs were actually consulted for this attempt;
retention of an additional original reference does not claim fresh consultation of it.
\begin{itemize}
\item \url{https://ems.press/content/serial-article-files/46961} (consulted 2026-10-11). 
\item \url{https://arxiv.org/html/2410.18661v2} (consulted 2026-10-11). 
\end{itemize}
\section{Model and contributor disclosure}
Codex GPT-6 completed written mathematical review of the full reduction and scope; this is model-assisted review, not independent human proof checking.
OpenAI gpt-6-astra, requested reasoning configuration ultra, generated this research
attempt on 2026-10-11. The exact serving revision is not exposed. Codex GPT-6
performed orchestration and repository-format review and suggested the material integer-to-cyclic reduction; its exact serving revision
is also not exposed. The mathematical route, boundary checks and final adversarial
self-audit followed contribution prompts A then B in
\texttt{docs/contribution\_prompts.md}, after reading \texttt{CONTRIBUTING.md}.
Python 3.14.2 assembled the standalone documents; any mathematical computation is
specified in Verification. No theorem prover was used. No human mathematical
authorship, substantial human proof correction, or independent human proof checking
is asserted. Model-assisted review, source retrieval, and structural checks do not
establish correctness or novelty. Full private reasoning transcripts are unavailable
as reproducibility artifacts; the complete arguments, inputs and outcomes actually
used are stated here or in the companion batch audit.
\section{Final status}
Solved as a complete complexity classification, conditional on the explicitly cited
known integer NP-completeness theorem. A polynomial algorithm exists exactly if
$\mathrm P=\mathrm{NP}$. No unconditional separation, novelty, or independently
verified proof of the external integer gadgets is asserted. The historical catalogue
source status remains unchanged.
\end{document}
